% Section 6.4, from lectures/L06/appendix/c_exercises.md.

\section{Övningar}
\label{c6:sec:exercises}\label{c6:app:c}

\begin{aside}
\textbf{Så kontrollerar du ditt arbete.} Ett sekvensnät kan inte kontrolleras rad för rad i en
sanningstabell, eftersom samma ingångar kan ge olika utgång beroende på vad som hänt innan. Gå i
stället igenom en \textbf{följd} av ingångar, ett steg i taget, och skriv upp tillståndet efter
varje steg. I CircuitVerse gör du samma sak: ändra en ingång, titta på utgången, anteckna, nästa.

Kontrolluppgiften är övning~\ref{c6:ex:4}. Lösningsförslagen finns i bilaga~\ref{app:answers}.
\end{aside}

Varje övning är märkt med sin sort.

\exercise{Kombinatoriskt eller sekventiellt?}{Förståelse}
\label{c6:ex:1}
Avgör för var och en av kretsarna nedan om den är kombinatorisk eller sekventiell, och motivera med
en mening.
\begin{parts}
  \item Tankens larm i kapitel~\ref{c4:ch}, som går när minst två av tre givare ger 1.
  \item Start- och stoppkretsen med självhållning i kapitel~\ref{c2:ch}.
  \item En avkodare som tänder rätt segment i en sifferdisplay för en BCD-siffra.
  \item Ett trafikljus som växlar mellan rött, gult och grönt.
  \item Ett kodlås som öppnar när siffrorna 3, 7 och 1 har tryckts i den ordningen.
  \item Bromsassistenten i kapitel~\ref{c2:ch}.
\end{parts}

\exercise{SR-låset}{Räkna för hand}
\label{c6:ex:2}
Ett SR-lås av två NOR-grindar (\secref{c6:a3}) börjar med \code{Q = 0}. Ingångarna ändras i den här
ordningen. Fyll i \code{Q} efter varje steg.

\begin{filltable}{@{}YYYY@{}}
  \thead{Steg} & \thead{S} & \thead{R} & \thead{Q} \\
  \midrule
  start & 0 & 0 & 0 \fillrule
  1 & 1 & 0 & \fillrule
  2 & 0 & 0 & \fillrule
  3 & 0 & 0 & \fillrule
  4 & 0 & 1 & \fillrule
  5 & 0 & 0 & \fillrule
  6 & 1 & 0 & \fillrule
  7 & 0 & 0 & \\
\end{filltable}

\begin{parts}
  \item Fyll i tabellen.
  \item Steg 2, 3, 5 och 7 har samma ingångar. Har de samma \code{Q}? Varför är det här skälet till
    att nätet kallas sekventiellt?
  \item Varför är kombinationen \code{S = R = 1} förbjuden i ett SR-lås av NOR-grindar?
\end{parts}

\exercise{Lås och vippa}{Räkna för hand}
\label{c6:ex:3}
Ett D-lås och en D-vippa får samma klocka och samma \code{D}. Låset är genomsläppligt medan
klockan är hög; vippan läser \code{D} vid varje stigande flank. Båda börjar med \code{Q = 0}.

\bookfigure{l06_ex_timing}{Klockan och \code{D} för övning~\ref{c6:ex:3}.}{c6:fig:ex_timing}

\begin{parts}
  \item Rita \code{Q} för låset.
  \item Rita \code{Q} för vippan.
  \item \code{D} har en kort puls mellan den andra och den tredje stigande flanken. Syns den på
    låsets utgång? På vippans? Förklara.
\end{parts}

\exercise{Kontroll: bygg ett SR-lås}{Kontroll}
\label{c6:ex:4}
\emph{Förutsäg, bygg i CircuitVerse, jämför.}
\begin{parts}
  \item \textbf{Förutsäg.} Använd din tabell från övning~\ref{c6:ex:2}.
  \item \textbf{Bygg} SR-låset av två NOR-grindar i CircuitVerse
    (\url{https://circuitverse.org/simulator}), med två ingångar \code{S} och \code{R} och två
    utgångar \code{Q} och \code{Q'}, korskopplade som i \secref{c6:a3}.
  \item \textbf{Kör följden} från övning~\ref{c6:ex:2}, ett steg i taget. Anteckna \code{Q} och
    \code{Q'} efter varje steg, och jämför med din förutsägelse.
  \item När du precis har byggt kretsen, innan du tryckt på något, vad visar \code{Q}? Är det samma
    sak varje gång du bygger om den? Vad säger det om ett minne som slås på?
  \item Ställ in \code{S = R = 1} och sedan båda till 0 samtidigt. Vad händer, och stämmer det med
    vad \secref{c6:a3} säger?
\end{parts}

\exercise{Klockan}{Räkna för hand}
\label{c6:ex:5}
\begin{parts}
  \item En Arduino Uno har klockfrekvensen 16~MHz. Hur lång är en klockcykel?
  \item Hur många klockcykler går det på en millisekund? På en sekund?
  \item Ett program utför 1000 instruktioner som tar en klockcykel var. Hur lång tid tar det?
  \item Samma program på en krets som går på 8~MHz. Hur lång tid då?
\end{parts}

\exercise{Räknare}{Räkna för hand}
\label{c6:ex:6}
\begin{parts}
  \item En 3-bitars räknare står på \code{000}. Vad står den på efter 11 stigande klockflanker?
  \item Hur många vippor behövs för en räknare som ska kunna räkna från 0 till 999?
  \item En 4-bitars räknare får en klocka på 16~Hz. Med vilken frekvens växlar den mest
    signifikanta biten, \code{Q3}?
  \item Ett 8-bitars register står på 255, och räknas upp ett steg. Vad står det på efter? Varför?
\end{parts}

\exercise{Skiftregister}{Räkna för hand}
\label{c6:ex:7}
Ett 4-bitars skiftregister som i \secref{c6:a8} börjar med \code{Q0 Q1 Q2 Q3 = 0000}.
\begin{parts}
  \item Bitarna 1, 1, 0 och 1 skickas in på \code{in}, en per klockflank. Skriv upp
    \code{Q0}--\code{Q3} efter varje flank.
  \item Läs \code{Q3 Q2 Q1 Q0} som ett binärt tal. Registret innehåller \code{0011} (talet 3), och
    ett steg till skiftas in med \code{in = 0}. Vilket tal står det nu? Vilken räkneoperation
    motsvarar ett steg?
\end{parts}

\exercise{Blockschemat}{Förståelse}
\label{c6:ex:8}
Vilket block i mikrodatorn (\secref{c6:b2}) beskrivs?
\begin{parts}
  \item Håller adressen till nästa instruktion.
  \item Adderar två tal, eller jämför dem.
  \item Håller programmet, även när strömmen är avslagen.
  \item Talar om att det senaste resultatet blev noll.
  \item Tolkar instruktionen och talar om för de andra blocken vad de ska göra.
  \item Ett register vars bitar är kopplade till kretsens stift.
  \item Håller variablerna medan programmet kör.
  \item Ger takten som alla vippor uppdateras i.
\end{parts}

\exercise{Instruktionscykeln}{Räkna för hand}
\label{c6:ex:9}
Programminnet innehåller:

\begin{narrowtable}{@{}cll@{}}
  \thead{Adress} & \thead{Instruktion} & \thead{Betyder} \\
  \midrule
  0 & \code{ldi r16, 10} & lägg 10 i r16 \\
  1 & \code{ldi r17, 20} & lägg 20 i r17 \\
  2 & \code{add r16, r17} & r16 = r16 + r17 \\
  3 & \code{rjmp 3} & hoppa till adress 3 \\
\end{narrowtable}

\begin{parts}
  \item Fyll i en tabell som den i \secref{c6:b4} för de sex första instruktionerna som utförs: PC
    före, instruktionen, vad som händer, r16 och r17 efter, PC efter.
  \item Vad gör programmet efter steg 4, och hur länge? Varför behövs en sådan instruktion sist i
    ett program för en mikrokontroller?
\end{parts}

\exercise{Tre minnen}{Förståelse}
\label{c6:ex:10}
\begin{parts}
  \item I vilket av ATmega328P:s tre minnen hamnar: programmet; en variabel som räknar hur många
    gånger en knapp tryckts; ett kalibreringsvärde som ska finnas kvar efter ett strömavbrott?
  \item Varför finns programmet kvar när strömmen slås av, men inte variablerna?
  \item Vad betyder det att ATmega328P har en Harvardarkitektur?
\end{parts}

\exercise{Bussar och minnesstorlek}{Räkna för hand}
\label{c6:ex:11}
\begin{parts}
  \item ATmega328P har 2~kB SRAM, alltså 2048 byte. Hur många adressbitar behövs för att peka ut
    varje byte?
  \item Flashminnet är 32~kB och varje instruktion är 16 bitar (2 byte). Hur många instruktioner
    ryms som mest, och hur många bitar behöver programräknaren?
  \item Hur många olika värden kan databussen i en 8-bitars dator bära?
\end{parts}

\exercise{Självhållning och SR-lås}{Konstruktion, Fördjupning}
\label{c6:ex:12}
Titta på självhållningskretsen i \secref{c6:a2}.
\begin{parts}
  \item Vilken knapp motsvarar SR-låsets \code{S}, och vilken \code{R}?
  \item Vad händer i reläkretsen om båda knapparna trycks samtidigt? Jämför med SR-låset av
    NOR-grindar.
  \item I en del maskiner är det tvärtom önskvärt att start vinner om båda trycks, kallat
    \textbf{tillslagsdominerande} självhållning. Rita om kretsen så att \code{S1} alltid kan dra
    \code{K1}, även när \code{S2} är nedtryckt. Varför är den varianten sällsynt för maskiner där
    stopp är en säkerhetsfunktion?
\end{parts}
